% Part: methods
% Chapter: induction
% Section: induction-on-N

\documentclass[../../../include/open-logic-section]{subfiles}

\begin{document}

\olfileid{mth}{ind}{inN}

\olsection{~$\Nat$ اُتے استقرا}

اپنی سبھ توں سوکھی صورت وچ استقرا سارے قدرتی عدداں لئی نتیجے ثابت کرن دا طریقہ اے۔ ایہ ایس گل نوں ورتدا اے کہ $0$ توں شروع کر کے وار وار~$1$ پائیے تاں آخر ہر قدرتی عدد تک پہنچ جاندے آں۔ سو ہر عدد لئی کوئی گل سچ ثابت کرن واسطے اسیں (۱) ثابت کر سکدے آں کہ ایہ $0$ لئی سچ اے، تے (۲) وکھا سکدے آں کہ جدوں وی ایہ کسے عدد~$n$ لئی سچ اے، تاں اگلے عدد~$n+1$ لئی وی سچ اے۔ جے اسیں «عدد~$n$ وچ خاصیت $P$ اے» نوں مختصر $P(n)$ لکھیے (تے «عدد~$k$ وچ خاصیت $P$ اے» نوں $P(k)$، وغیرہ)، تاں ایہ ثابت کرن لئی کہ $P(n)$ سارے $n \in \Nat$ لئی، استقرائی ثبوت وچ ایہ ہُندا اے:
\begin{enumerate}
\item $P(0)$ دا ثبوت، تے
\item ایہدا ثبوت کہ کسے وی~$k$ لئی، جے $P(k)$ تاں $P(k+1)$۔
\end{enumerate}
گل پوری صاف کرن لئی منّ لو ساڈے کول (۱) تے~(۲) دوویں نیں۔ (۱) دسدا اے کہ $P(0)$ سچ اے۔ جے~(۲) وی اے، تاں سانوں خاص کر کے پتا اے کہ جے $P(0)$ تاں $P(0+1)$، یعنی $P(1)$۔ ایہ عام بیان «کسے وی~$k$ لئی، جے $P(k)$ تاں $P(k+1)$» توں~$0$ نوں~$k$ دی تھاں رکھ کے نکلدا اے۔ سو شرطیہ تے اوہدے مقدم توں نتیجہ کڈھن دے قاعدے نال سانوں~$P(1)$ ملدا اے۔ (۲) توں فیر، ہُن $1$ نوں~$n$ دی تھاں رکھ کے، سانوں ملدا اے کہ جے $P(1)$ تاں~$P(2)$۔ کیوں جے اسیں ہُن ای~$P(1)$ ثابت کیتا اے، اوسے قاعدے نال سانوں~$P(2)$ ملدا اے۔ تے ایس طرح اگے۔ کسے وی عدد~$n$ لئی، ایہ کم $n$~واری کر کے آخر~$P(n)$ تک پہنچدے آں۔ سو (۱) تے~(۲) رَل کے~$P(n)$ ثابت کردے نیں کسے وی $n
\in \Nat$ لئی۔

اک مثال ویکھیے۔ منّ لو اسیں پتا لاؤنا چاہندے آں کہ $n$ پاسے سُٹ کے کِنّے وکھرے جوڑ کڈھ سکدے آں۔ بھاویں ایہ ذرا اوپرا لگے، آؤ $0$ پاسیاں توں شروع کریے۔ جے تہاڈے کول کوئی پاسا نہیں، تاں بس اک ممکن جوڑ «سُٹ» سکدے او۔ کوئی نقطہ نہیں، جیہدا جوڑ~$0$ اے۔ سو وکھرے ممکن حاصلاں دی گنتی~$1$ اے۔ جے اکّو پاسا اے، یعنی $n=1$، تاں چھ ممکن قدراں نیں، $1$ توں~$6$ تک۔ دو پاسیاں نال $2$ توں $12$ تک کوئی وی جوڑ سُٹ سکدے آں، یعنی $11$ صورتاں۔ تِن پاسیاں نال $3$ توں $18$ تک کوئی وی عدد سُٹ سکدے آں، یعنی $16$ وکھریاں صورتاں۔ $1$، $6$، $11$، $16$، اک ترتیب لگدی اے۔ شاید جواب $5n+1$ اے؟ بے شک $5n+1$ سبھ توں ودھ ممکن گنتی اے، کیوں جے بس $5n+1$ عدد نیں $n$ توں، جیہڑی $n$ پاسیاں نال سبھ توں گھٹ قدر اے (سارے $1$)، $6n$ تک، جیہڑی سبھ توں ودھ قدر اے (سارے $6$)۔

\begin{thm}
  $n$ پاسیاں نال ساریاں $5n+1$ ممکن قدراں سُٹ سکدے آں جیہڑیاں $n$ تے $6n$ دے وچکار نیں۔
\end{thm}

\begin{proof}
$P(n)$ ایہ دعویٰ ہووے، «$n$ تے~$6n$ دے وچکار کوئی وی عدد $n$~پاسیاں نال سُٹ سکدے آں»۔ استقرا ورتن لئی ایہ ثابت کردے آں:
\begin{enumerate}
\item \emph{استقرا دی بنیاد} $P(1)$، یعنی اکّو پاسے نال $1$ تے $6$ دے وچکار کوئی وی عدد سُٹ سکدے او۔
\item \emph{استقرا دا قدم}، سارے $k$ لئی، جے $P(k)$ تاں~$P(k+1)$۔
\end{enumerate}

(۱) اک $6$ رُخاں والا پاسا ویکھ کے ثابت ہُندا اے۔ اوہدے سارے 6 رُخ نیں، تے $1$ توں~$6$ تک ہر عدد کسے اک رُخ اُتے اے۔ سو اکّو پاسے نال $1$ تے~$6$ دے وچکار کوئی وی عدد سُٹ سکدے آں۔

(۲) ثابت کرن لئی اسیں شرطیے دا مقدم منّدے آں، یعنی $P(k)$۔ ایس منّی ہوئی گل نوں \emph{استقرائی مفروضہ} آکھدے نیں۔ اسیں ایہنوں~$P(k+1)$ ثابت کرن لئی ورتدے آں۔ اوکھا حصہ ایہ لبھنا اے کہ $k+1$ پاسیاں دے سُٹّن دیاں ممکن قدراں نوں، $k$~پاسیاں دے سُٹّن تے وادھو $k+1$ ویں پاسے دے سُٹّن دیاں قدراں نال کیویں جوڑیے۔ جے استقرائی مفروضہ ورتنا اے تاں ایہ کم کرنا پئے گا۔

استقرائی مفروضہ آکھدا اے کہ $k$ تے~$6k$ دے وچکار کوئی وی عدد $k$~پاسیاں نال کڈھ سکدے آں۔ جے~$1$ اپنے $(k+1)$ ویں پاسے نال سُٹیے، تاں کُل جوڑ وچ $1$ ودھ جاندا اے۔ سو $k+1$ توں $6k+1$ تک کوئی وی قدر سُٹ سکدے آں، $k$~پاسے سُٹ کے تے فیر~$1$ نوں $(k+1)$ ویں پاسے نال سُٹ کے۔ کیہ رہ گیا؟ $6k+2$ توں $6k+6$ تک دیاں قدراں۔ ایہناں لئی $k$ واری $6$ سُٹو تے فیر $2$ تے $6$ دے وچکار کوئی عدد اپنے $(k+1)$ ویں پاسے نال سُٹو۔ رَل کے ایہدا مطلب اے کہ $k+1$ پاسیاں نال $k+1$ تے $6(k+1)$ دے وچکار کوئی وی عدد سُٹ سکدے آں، یعنی اسیں~$P(k+1)$ مفروضے~$P(k)$، یعنی استقرائی مفروضے، نال ثابت کر دتا اے۔
\end{proof}

اسیں استقرا اکثر اوہدوں ورتدے آں جدوں شےواں (عدد، سیٹ، وغیرہ) دی کسے لڑی بارے کوئی گل ثابت کرنی ہووے تے لڑی آپ «استقرائی طریقے نال» تعریف ہوئی ہووے، یعنی $(n+1)$ ویں شے دی تعریف $n$ ویں شے نال دتی ہووے۔ مثلاً قدرتی عدداں دا جوڑ~$s_n$،~$n$ تک، ایویں تعریف کر سکدے آں:
\begin{align*}
  s_0 & = 0\\
  s_{n+1} & = s_n + (n+1)
\end{align*}
ایس تعریف توں ملدا اے:
\begin{align*}
  s_0 & = 0,\\
  s_1 & = s_0 + 1 && = 1,\\
  s_2 & = s_1 + 2 && = 1 + 2 = 3\\
  s_3 & = s_2 + 3 && = 1 + 2 + 3 = 6, \text{ وغیرہ۔}
\end{align*}
ہُن اسیں استقرا نال ثابت کر سکدے آں کہ $s_n = n(n+1)/2$۔

\begin{prop}
  $s_n = n(n+1)/2$.
\end{prop}

\begin{proof}
  سانوں ثابت کرنا اے (۱) کہ $s_0 = 0\cdot(0 + 1)/2$، تے (۲) جے $s_k =
  k(k+1)/2$ تاں $s_{k+1} = (k+1)(k+2)/2$۔ (۱) صاف اے۔ (۲) ثابت کرن لئی اسیں استقرائی مفروضہ منّدے آں، $s_k = k(k+1)/2$۔ ایہ ورت کے سانوں وکھاؤنا اے کہ $s_{k+1} = (k+1)(k+2)/2$۔

  $s_{k+1}$ کیہ اے؟ تعریف مطابق $s_{k+1} = s_k + (k+1)$۔ استقرائی مفروضے مطابق $s_k = k(k+1)/2$۔ ایہ پچھلی مساوات وچ پا سکدے آں، تے فیر بس کسراں دا تھوڑا حساب چاہیدا اے:
  \begin{align*}
    s_{k+1} & = \frac{k(k+1)}{2} + (k+1) = {}\\
    & = \frac{k(k+1)}{2} + \frac{2(k+1)}{2} = {}\\
    & = \frac{k(k+1) + 2(k+1)}{2} = {}\\
    & = \frac{(k+2)(k+1)}{2}.
  \end{align*}
\end{proof}

ایتھے ضروری گل ایہ اے کہ جے تسی استقرائی طریقے نال تعریف ہوئی کسے لڑی $a_n$ بارے کوئی گل ثابت کردے او، تاں استقرا ورتنا سِدھا طریقہ اے۔ جے ایہ گل سِدھی نہ دِسے (جیویں پاسے سُٹّن دیاں ممکن صورتاں وچ)، تاں وی استقرا ورت سکدے او جے~$k+1$ والی صورت نوں~$k$ والی صورت نال کسے طرح جوڑ سکو۔

\end{document}
